\section{An evaluated Hecke separator and the missing integral extension data}
\label{sec:kt-hecke-normal}

The complementary stable factors computed earlier do give an
integral, all-prime statement about the localized Hecke algebra.
This is an arithmetic consequence of an evaluated unit.  The
following audit then explains precisely why it does not yet give
a smaller strict Selmer bound.

Write $A_q=q^{e-1}T_q$ and $A_q^{\rm opp}=q^{e-1}T_q^{\rm opp}$
for the two integral scaled split operators, and put
$t_q=1+q^e+q^{2e}$.  Let $M^{\rm st}_e$ be the stable compact
space, with its proved decomposition into the constructed
symmetric-square subspace and its complementary stable subspace.
Write $R_{e,7}(X)$ for the characteristic polynomial of $A_7$
on the latter.  Thus its degrees at $e=26,28,30$ are
$72,92,112$.

\begin{proposition}[KT-STABLE-HECKE-SELFDUAL-LOCALIZATION]
For $e\in\{26,28,30\}$ put, respectively,
\[
 M_e=697446897311,\quad323026426255399,\quad
      236017358500033404991.
                                                               \tag{KH.1}
\]
Let $\ell$ be any prime divisor of $M_e$.  On the entire stable
space one has, for every split $q\ne\ell$,
\[
 R_{e,7}(A_7)(A_q-A_q^{\rm opp})=0.                         \tag{KH.2}
\]
Choose a common integral $\mathbf Z_\ell$-lattice and let
$\mathbb T^{\rm st}$ be the algebra generated by all the displayed
split operators.  At the Eisenstein residual maximal ideal
\[
 \mathfrak m=(\ell,A_q-t_q,A_q^{\rm opp}-t_q:q\ne\ell
                   \text{ split}),                        \tag{KH.3}
\]
one has the integral identities
\[
         A_q=A_q^{\rm opp}\quad\hbox{in }
                       \mathbb T^{\rm st}_{\mathfrak m}.  \tag{KH.4}
\]
The involution interchanging the two operators is the identity
on this localized order.  Its characteristic-zero cuspidal
systems all belong to the constructed symmetric-square subspace.

The same assertion holds for the augmented order generated by
the pairs $(t_q,A_q)$ and $(t_q,A_q^{\rm opp})$ inside
$\mathbf Z_\ell\oplus\mathbb T^{\rm st}$, localized at its
augmentation modulo $\ell$.  Here the first factor is a formal
Eisenstein character; it is not asserted to be a compact
cuspidal representation.
\end{proposition}
\begin{proof}
All the split Hecke operators commute.  On the complementary
subspace Cayley--Hamilton makes $R_{e,7}(A_7)$ zero.  On the
constructed symmetric-square subspace the associated adjoint
representation is self-dual, and the two normalized minuscule
eigenvalues agree at every split prime.  Semisimplicity therefore
makes $A_q-A_q^{\rm opp}$ zero there.  These facts prove (KH.2)
over characteristic zero, and hence on the integral lattice.
No integral direct-sum decomposition of that lattice is needed.

The integral harmonic lattice at $\ell>3$ supplies a common
Hecke-stable lattice; intersect it with the rational stable
subspace.  The involution given by interchanging the harmonic
variables preserves it after replacing the lattice by its sum
with its image, if necessary.  KT-U3-DUAL-MATRIX identifies its
conjugation action on the two split operators.  The algebra
$\mathbb T^{\rm st}$ is a finite torsion-free
$\mathbf Z_\ell$-module, being a submodule of the finite matrix
algebra.  The earlier Fretwell--Roberts construction provides a
cuspidal eigencharacter with the values (KH.3), so this ideal is
proper and maximal.

The exact calculations give
$\gcd(M_e,R_{e,7}(t_7))=1$.  Consequently the image of
$R_{e,7}(A_7)$ modulo $\mathfrak m$ is nonzero and this element
is a unit in the localized algebra.  Multiplying (KH.2) by its
inverse proves (KH.4).  The generators are all fixed, so the
involution is the identity.  A characteristic-zero factor from
the complement cannot survive this localization, since on that
factor $R_{e,7}(A_7)=0$.  This also proves the generic-support
assertion without an assumption of distinct eigenvalues at seven.

On the added Eisenstein coordinate the second factor in (KH.2)
is zero.  The same relation and unit argument therefore apply
to the augmented order, which is again finite and torsion-free.
Its first-coordinate projection gives the specified augmentation.
\end{proof}

The new small certificate \texttt{hecke\_normal\_audit.py}
evaluates the previously reviewed exact coefficient arrays
modulo $M_e$ and saves explicit Bezout identities.  At seven
the reduced values and their inverses are
\[
\begin{array}{c|r|r}
e&R_{e,7}(t_7)\bmod M_e&R_{e,7}(t_7)^{-1}\bmod M_e\\\hline
26&429636065482&455711584154\\
28&242701611239023&44010470192816\\
30&53393982565377757302&202700793921978609871
\end{array}                                                   \tag{KH.5}
\]
The analogous three certificates at thirteen are also saved.
No primality assumption about any $M_e$ is used.  These six unit
certificates strengthen the finite-prime comparison to (KH.4),
but they are not a new Selmer vanishing certificate.

\subsection{Why a trace comparison does not detect the required tangent}

\begin{lemma}[KT-TRACE-INVISIBLE-SELMER-TANGENT]
Let $S,x_0,C$ be the actual normalized source ring, its diagonal
augmentation, and the odd weight $-1$ cotangent from
KT-ODD-CONORMAL-SELMER-EQUIVALENCE.  For every nonzero $b\in O$,
each element of $\operatorname{Hom}_O(C,O/b)$ induces zero
first-order variation of every characteristic coefficient of
the universal representation restricted to $G_E$.  In particular
it is zero on every trace coordinate and on every polynomial
in such coordinates.
\end{lemma}
\begin{proof}
The proved inverse map identifies the corresponding lift over
$O\oplus\epsilon O/b$, $\epsilon^2=0$, as
\[
 r_\epsilon(g)=(1+\epsilon t(g)Y)r_0(g),\qquad
 Y=E_{21}-E_{32},\qquad g\in G_E,                           \tag{KH.6}
\]
with the prescribed normalized strict Selmer cocycle $t$.
The matrix $r_0(g)$ is diagonal, so $r_\epsilon(g)$ is lower
triangular with exactly the same diagonal entries.  Its
characteristic polynomial is therefore equal to that of
$r_0(g)$, as an exact equality over this square-zero coefficient
ring.  This includes traces of every group element, hence also
traces of all products of group matrices.  Equivalently the
displayed first derivative is
\[
             \operatorname{tr}(t(g)Y r_0(g))=0.            \tag{KH.7}
\]
The equality is integral and uses no semisimplification of the
first-order representation.  The product rule proves the claim
for polynomial trace coordinates.  The weight description gives
the same conclusion: such coordinates have torus weight zero,
while this derivation is supported on weight $-1$.
\end{proof}

Consequently restricting these derivations to the trace
subalgebra gives the zero map.  Showing this restriction to be
injective would already prove that the original strict Selmer
module vanishes.  The identity involution in (KH.4) supplies no
such injectivity.  Even an identification of an integral
pseudorepresentation algebra with the displayed Hecke order
would need an additional theorem controlling the off-diagonal
extension modules.  A generalized matrix algebra comparison,
or an equivalent comparison of integral lattices and their
reducibility pairings, could in principle provide those data;
neither the characteristic-zero source coverage nor the actual
point construction is that comparison.  The original classical
rank-two branch is also absent from the stable Hecke algebra and
cannot be removed by (KH.4).

\subsection{The trace pairing and Euler characteristic do not supply a unit}

There is a further concrete obstruction to treating an
irreducible adjoint point as integrally nondegenerate.  It
explains where denominators enter an attempt to recover matrix
entries from traces.

\begin{lemma}[KT-RESIDUAL-TRACE-DISCRIMINANT]
Let $r:G_E\to\mathrm{GL}_3(O')$ be an actual adjoint lift with
the split diagonal residual lattice constructed in
KT-ADJOINT-ACTUAL-HOST-POINT.  Suppose, as holds for that
non-CM lift, that $r$ is absolutely irreducible over the fraction
field $K'$ of the coefficient DVR $O'$.  Let $\mathcal A$ be
the $O'$-span of the matrices $r(g)$.  Then $\mathcal A$ is a
full lattice in $M_3(K')$, and the discriminant of its trace
pairing satisfies
\[
 v_{O'}\operatorname{disc}(\mathcal A)
       =2\operatorname{length}_{O'}(M_3(O')/\mathcal A)
       \geq12.                                             \tag{KH.8}
\]
In particular this trace discriminant is not a unit.
\end{lemma}
\begin{proof}
Absolute irreducibility gives the full matrix algebra as the
$K'$-span of the group, so $\mathcal A$ has rank nine.  Its
reduction inside $M_3(k')$ is contained in the three-dimensional
diagonal subspace.  At least six elementary divisors of its
inclusion in $M_3(O')$ are therefore divisible by the uniformizer,
giving the lower bound on the length.  The matrix-unit basis of
$M_3(O')$ has unimodular pairing $(X,Y)\mapsto\operatorname{tr}(XY)$.
Changing from that basis to a basis of $\mathcal A$ multiplies
the Gram determinant by the square of the inclusion determinant.
Smith normal form proves (KH.8).
\end{proof}

This is a statement about the actual split residual adjoint
lattice, not an evaluation of the odd normal Jacobian.  It
does not imply that the strict Selmer group is nonzero; it
only prevents replacing integral recovery of matrix entries by
the characteristic-zero density theorem without its denominators.

The local and global cohomological dimension formulas used
earlier also do not determine those denominators.  For example,
the two-term complex
$[O\xrightarrow{\varpi^n}O]$ has zero rational cohomology and
the same module ranks for every $n\geq1$, but its integral
cokernel has length $n$.  A perfect finite pairing alone leaves
the same ambiguity: $(x,y)\mapsto xy/\varpi^n\bmod O$ is a
perfect pairing on $O/\varpi^n$ for every $n$.  These are algebraic
scope examples, not counterexamples to the source hypotheses.
For the rational character $\chi^{-e}$ the Tate-dual coefficient
is $\chi^{e+1}$; the duality operation does not identify the
required integral group with a group already proved zero here.

At an adjoint branch the already proved odd normal coefficient
is $\eta\otimes(\mathbf1\oplus\operatorname{Sym}^4\rho
\otimes(\det\rho)^{-2})$.  Characteristic-zero rigidity proves
full generic rank of its normal Jacobian, but no normalized
arithmetic determinant or integral special-value comparison for
that Jacobian has been supplied.  Nor has its base change to the
Eisenstein augmentation been bounded independently of the
original strict Selmer module.  Thus the evaluated unit (KH.5)
does prove the new Hecke identity, while a strict Selmer upper
bound would still require the missing integral extension-module
comparison and an evaluated annihilator or normal determinant.
No stronger arithmetic vanishing conclusion is claimed.

\paragraph{A published comparison route and its failed hypothesis here.}
A targeted primary-source check finds a direct model for the
missing comparison in Hida--Tilouine, \emph{Symmetric Power
Congruence Ideals and Selmer Groups}, author manuscript dated
7 October 2018, pp.~2--4: residual irreducibility is imposed
on p.~2, hypothesis $(*)$ on p.~3 requires residual image
containing $\mathrm{SL}_2(k')$, and Theorems~1.1 and~1.3 use
these hypotheses for the deformation/congruence-ideal comparison.
The source is \url{https://www.math.ucla.edu/~hida/321.pdf}.
Zhang's \emph{Selmer Groups of Symmetric Powers of Ordinary
Modular Galois Representations}, arXiv:1802.08329v1, printed
p.~8, imposes the residual condition
$\operatorname{Sym}^{n-1}\mathrm{SL}_2(\mathbf F')
\subset\operatorname{im}\overline\rho$ in Condition~2.7;
Theorem~2.11 assumes this condition in its $R=T$ conclusion.
The exact version is
\url{https://arxiv.org/pdf/1802.08329v1}.

For the present classical forms,
$\overline\rho_f^{\rm ss}=\mathbf1\oplus\eta\overline\chi^e$.
Every residual lattice is therefore reducible, including a
possible nonsplit extension; its symmetric powers preserve
a nontrivial flag.  Scalar twisting or restricting the Galois
group preserves that flag.  Thus the stated residual-image
hypotheses fail for these branches.  These papers supply no
applicable integral comparison here on their stated hypotheses.
Only those hypotheses and relevant theorem statements were
used in this exclusion audit; their full comparison proofs
have not been imported as new inputs.
